% Part: first-order-logic
% Chapter: syntax
% Section: intro-syntax

\documentclass[../../../include/open-logic-section]{subfiles}

\begin{document}

\olfileid{fol}{syn}{itx}

\olsection{مقدمه}

برای بسط نظریه و فرامنطق منطق مرتبهٔ اول، نخست باید نحو و معناشناسی
عبارت‌های آن را تعریف کنیم. عبارت‌های منطق مرتبهٔ اول ترم‌ها و
!!{formula}s هستند. ترم‌ها از !!{variable}s، !!{constant}s و
!!{function}s تشکیل می‌شوند. !!^{formula}s نیز به‌نوبهٔ خود از
!!{predicate}s همراه با ترم‌ها ساخته می‌شوند (این‌ها ساده‌ترین
!!{formula}s، یعنی فرمول‌های «اتمی»، را تشکیل می‌دهند)؛ سپس می‌توانیم
از !!{formula}s اتمی با استفاده از پیوندگرهای منطقی و سورها فرمول‌های
پیچیده‌تر بسازیم. برای صورت‌بندی قواعد تشکیل راه‌های گوناگونی وجود دارد
و ما تنها یکی از راه‌های ممکن را ارائه می‌کنیم. دستگاه‌های دیگر نمادهای
متفاوتی برمی‌گزینند، مجموعه‌های متفاوتی از پیوندگرها را به‌عنوان نمادهای
اولیه انتخاب می‌کنند و پرانتزها را به شیوه‌ای متفاوت به کار می‌برند
(یا حتی اصلاً به کار نمی‌برند، مانند آنچه در نمادگذاری موسوم به لهستانی
رخ می‌دهد). بااین‌حال، وجه مشترک همهٔ رویکردها آن است که قواعد تشکیل،
مجموعهٔ ترم‌ها و !!{formula}s را \emph{به‌طور استقرایی} تعریف می‌کنند.
اگر این کار درست انجام شود، هر عبارت بنا بر قواعد تشکیل اساساً فقط به
یک شیوه می‌تواند به دست آید. تعریف استقرایی‌ای که عبارت‌هایی
\emph{با خوانش یکتا} به دست می‌دهد، بدین معناست که می‌توانیم با همان
روش---یعنی تعریف استقرایی---به این عبارت‌ها معنا بدهیم.

\end{document}
